% master_index.tex
\documentclass[12pt]{article}
\usepackage[utf8]{inputenc}
\usepackage[paper=a4paper,margin=2.5cm]{geometry}
\usepackage{hyperref}
\begin{document}

\title{אינדקס של אינדקסים לחבילת תבניות המחקר}
\author{מחולל פרומטים}
\date{\today}
\maketitle

\section*{מטרה}
קובץ זה מרכז את כל האינדקסים של הקבצים בחבילה ומספק קישורים ותמצית לכל אינדקס.

\section*{רשימת אינדקסים}
\begin{itemize}
  \item \textbf{index\_prompt\_generator.tex} -- תיאור מבנה הפרומט והסעיפים.
  \item \textbf{index\_parameters.tex} -- תיאור כל פרמטר והמלצות ערכים.
  \item \textbf{index\_extraction\_list.tex} -- שדות חילוץ ותבניות פלט.
  \item \textbf{index\_research\_template.tex} -- מבנה מסמך מחקר לדוגמה.
  \item \textbf{index\_README.tex} -- הוראות שימוש ותזכירים.
\end{itemize}

\section*{כיצד להשתמש}
פתחו את האינדקס המתאים מהסעיף למעלה כדי לעבור ישירות לתיאור המפורט של הקובץ ולרשימת הסעיפים.

\end{document}
